\origpage{153}{GRAPH MINORS, X}
{\centering
{\Large Graph Minors. X. Obstructions to Tree-Decomposition\par}\medskip

Neil Robertson$^*$\par
\textit{Department of Mathematics, Ohio State University, 231 West 18th Avenue, Columbus, Ohio 43210}\par
and\par
P. D. Seymour\par
\textit{Bellcore, 445 South Street, Morristown, New Jersey 07960}\par
\textit{Communicated by the Managing Editors}\par
Received May 10, 1988\par
}
\begin{quote}
Roughly, a graph has small ``tree-width'' if it can be constructed by piecing small graphs together in a tree structure. Here we study the obstructions to the existence of such a tree structure. We find, for instance:

\hspace*{12pt}(i) a minimax formula relating tree-width with the largest such obstructions

\hspace*{12pt}(ii) an association between such obstructions and large grid minors of the graph

\hspace*{12pt}(iii) a ``tree-decomposition'' of the graph into pieces corresponding with the obstructions.

These results will be of use in later papers.

\end{quote}

\sectionhead{1. Tangles}

Graphs in this paper are finite and undirected and may have loops or multiple edges. The vertex- and edge-sets of a graph $G$ are denoted by $V(G)$ and $E(G)$. If $G_1 = (V_1, E_1)$, $G_2 = (V_2, E_2)$ are subgraphs of a graph $G$, we denote the graphs $(V_1 \cap V_2, E_1 \cap E_2)$ and $(V_1 \cup V_2, E_1 \cup E_2)$ by $G_1 \cap G_2$ and $G_1 \cup G_2$, respectively. A \textit{separation} of a graph $G$ is a pair $(G_1, G_2)$ of subgraphs with $G_1 \cup G_2 = G$ and $E(G_1 \cap G_2) = \emptyset$, and the \textit{order} of this separation is $|V(G_1 \cap G_2)|$.

It sometimes happens with a graph $G$ that for each separation $(G_1, G_2)$ of $G$ of low order, we may view one of $G_1, G_2$ as the ``main part'' of $G$, in%
\origfootnote{This work was performed under a consulting agreement with Bellcore.}
